\section{The ideal potential and its profiles}\label{sec:profiles}

We now choose the trial vector for the capacity comparison.
The free potential includes every positive number of periods, so its
capacity already contains the full ideal-magnon density.  A spatial
cutoff makes it a finite-volume test vector without truncating those
periods.  We also obtain the spatial and temporal estimates needed to
compare its evolution with the walk among holes.

Fix $S$, put $\ell=2S$, and use the scales in \eqref{eq:scales}.
Let $E=D_0\cup F$, where $F$ is a finite deterministic set of
slot-time pins, and let $i=(x,a)$ be an unpinned post-seam slot at time
zero.  Throughout the comparison assume that $B_{50R_*}(x)$ lies in
the interior of the cube and that no point of $E$ lies above this
ball.  The picture law is $\PP_P$, where $P=E\cup\{i\}$.
At a fixed exposure let $U$ be the forward one-period kernel killed
on $E$, retaining its zero rows and columns on killed cut slots.
Set $A=I-U$ and use the energy $e_y$ of
Section~\ref{sec:capacity}.

\subsection{The all-turn free potential}

On $\Z^3\times\{1,\ldots,\ell\}$, let $\mathscr L$ have rate
$1/2$ on every slot edge, and let $T_s=\exp(s\mathscr L)$.
Let $J$ be the projection that averages the slots at each site.
For a slot-uniform vector the site jump rate is $\ell/2=S$ to each
neighbor.  Denote this site semigroup by $P_s$, and its kernel from
the origin by $p_s(z)$.  Its Fourier multiplier is
\begin{equation}\label{eq:freemultiplier}
 \widehat p_s(k)
 =\exp\!\left[-2sS\sum_{j=1}^3(1-\cos k_j)\right],
 \qquad k\in[-\pi,\pi]^3.
\end{equation}
If a vector has zero sum on the slots at every site, then
$\mathscr L$ acts on it as $-3\ell I$.  Therefore
\begin{equation}\label{eq:slotdecomposition}
 T_s=P_sJ+e^{-3\ell s}(I-J).
\end{equation}
The free seam is $J$.  It commutes with $T_s$, so the ideal full-period
operator, in either direction, is
\[
 U_0=JT_\beta=T_\beta J=P_\beta J,
 \qquad
 (U_0^n)_{(z,b),(x,a)}=\frac1\ell p_{n\beta}(z-x)
 \quad(n\ge1).
\]
In particular the slot modes removed by the seam contribute nothing
to a positive full turn.

With $\eta=\ind_{\{i\}}$ on the full slot lattice, define the
slot-uniform function
\begin{equation}\label{eq:freepotential}
 v(z)=\frac1\ell\sum_{n\ge1}p_{n\beta}(z-x),
 \qquad c=\frac1{1+v(x)}.
\end{equation}
The estimates below show that the sum is finite and
$v(x)=O_S(\beta^{-3/2})$.  Positivity and the semigroup property give
\begin{equation}\label{eq:freepotentialidentity}
 U_0(\eta+v)=v.
\end{equation}
Choose $\chi:\R^3\to[0,1]$ smooth, equal to one when
$|z-x|_\infty\le R/2$ and zero when $|z-x|_\infty\ge R$, with
first and second derivatives bounded by $C/R$ and $C/R^2$.
Use its restriction to sites as a slot-uniform
multiplier.  The full trial vector is
\begin{equation}\label{eq:trialvector}
 h=c\chi(\eta+v),
\end{equation}
and the same symbol denotes its restriction to the cut holes.
It vanishes on every killed cut slot and satisfies $h(i)=1$.
Moreover $p_s(z)\le p_s(0)$, by the Fourier representation, so
$0\le h\le1$.

\subsection{Forward and reverse profiles}

Put $g=c(\eta+v)$ on the full slot lattice.  The forward and reverse
profiles, before restriction to holes, are
\begin{align}
 a_s&=\chi T_s g\quad(0\le s<\beta),
 &a_\beta&=\chi JT_\beta g,                         \label{eq:forwardprofile}\\
 b_0&=\chi g,\qquad b_{0+}=\chi Jg,
 &b_s&=\chi T_sJg\quad(0<s\le\beta).                 \label{eq:reverseprofile}
\end{align}
Here $s$ is elapsed time in the indicated direction.  The forward
traversal contains the seam at its end.  The reverse traversal contains
it at its beginning, with $b_0$ its incoming value and $b_{0+}$ its
outgoing value.  Each traversal includes the seam exactly once.
At every slice, the same symbols also denote the restrictions to the
holes on the specified side of any instantaneous operation; ``full
profile values'' will always refer to the formulas before restriction.

Multiplication by $\chi$ commutes with $J$.  Thus the full profiles
undergo the ideal seam operation, and \eqref{eq:freepotentialidentity}
gives the endpoint identities
\begin{equation}\label{eq:profileendpoints}
 a_0=b_0=h,\qquad a_\beta=b_\beta=h-c\eta.
\end{equation}
The smooth part of either profile is the contribution of $cv$;
the point part is the contribution of $c\eta$.  This decomposition
also applies to full profile values.

For $r_z=|z-x|_\infty$ and $d\ge1$, introduce
\begin{equation}\label{eq:weights}
 \begin{split}
 G(z)&=\beta^{-1}(D_*+r_z)^{-1},\qquad
 W(z)=\beta^{-1}(D_*+r_z)^{-2},\\
 P_d(z)&=d^{-3}(1+r_z/d)^{-12},\qquad
 W_d(z)=d^{-4}(1+r_z/d)^{-12}.
 \end{split}
\end{equation}
For a slot function $u$, a bounded-distance difference at $z$ means
$u(z,a)-u(z',a')$ with $|z-z'|_\infty$ at most a fixed constant;
the constants below may depend on that constant.  This includes
differences between two slots at the same site.

\begin{lemma}[Bounds for the free profiles]\label{lem:freebounds}
For fixed $S$ and all sufficiently large $\beta$, the following
bounds hold uniformly in the direction, elapsed time, site, and slot.
For the full smooth profile $u_s$, its absolute value,
bounded-distance differences, and open-interval time derivative are
bounded, respectively, by
\begin{equation}\label{eq:smoothbounds}
 C G(z),\qquad C W(z),\qquad C W(z)/D_*.
\end{equation}
For the full point profile $w_s$, put $d=\sqrt{1+s}$.
The corresponding bounds are
\begin{equation}\label{eq:pointbounds}
 C P_d(z),\qquad C W_d(z),\qquad C W_d(z)/d.
\end{equation}
The smooth profile is slot-uniform and does not change at the seam.
The slot differences and seam changes of the point profile satisfy
the difference bound in \eqref{eq:pointbounds}.
For each fixed $a>0$, when $s\ge a\beta$ the bounds
\eqref{eq:smoothbounds} also hold for $w_s$, with a constant depending
on $a$.
\end{lemma}

\begin{proof}
We first record the heat-kernel estimates before applying the cutoff.
For every fixed integer $m\ge0$ and $j=0,1,2$, Fourier inversion gives
\begin{equation}\label{eq:heatderivative}
 |\nabla^j p_t(z)|
 \le C_{m,j}(1+t)^{-(3+j)/2}
       \left(1+\frac{|z|_\infty}{\sqrt{1+t}}\right)^{-m}.
\end{equation}
Here $j$ nearest-neighbor differences may be taken in arbitrary
coordinate directions.  One time derivative satisfies the same bound
with $j=2$.
To verify the estimates, write
$\omega(k)=2S\sum_j(1-\cos k_j)$ and use
$\omega(k)\ge c_S|k|^2$ on $[-\pi,\pi]^3$.
A $j$-fold difference inserts a multiplier vanishing to order $j$
at zero; a time derivative inserts $-\omega$, which vanishes to
order two.  The integral without frequency differentiation is thus
$O((1+t)^{-(3+j)/2})$.
For the spatial weight, integrate by parts $m$ times in a coordinate
for which $|z_j|=|z|_\infty$.  The integral of the absolute value
of the differentiated multiplier is at most
$C_{m,j}(1+t)^{(m-3-j)/2}$: near zero this follows by rescaling
$k$ by $(1+t)^{-1/2}$, using
$\nabla\omega=O(|k|)$ and bounded higher derivatives; away from
zero the exponential dominates every power of $t$.
The periodic integrands have no boundary terms.  Combining this
bound with the bound without integration by parts proves
\eqref{eq:heatderivative}.  The case $0\le t\le1$ follows from
uniform bounds on the same derivatives.

The slot-uniform part of the smooth evolution is a sum of site kernels
at times $n\beta+s$, $n\ge1$.  For $j=0,1,2$, choose $m>j+1$ in
\eqref{eq:heatderivative}.  Splitting the time sum at $n\beta+s$
of order $r_z^2$, or comparing it with its integral on dyadic
intervals, gives
\begin{equation}\label{eq:heatsum}
 \sum_{n\ge1}(n\beta+s)^{-(3+j)/2}
     \left(1+\frac{r_z}{\sqrt{n\beta+s}}\right)^{-m}
 \le \frac{C}{\beta}(D_*+r_z)^{-1-j}.
\end{equation}
For completeness, when $r_z\le D_*$ the sum without its spatial
factor is $O(\beta^{-(3+j)/2})$.  When $r_z>D_*$, the terms with
$n\beta+s\ge r_z^2$ sum to $C\beta^{-1}r_z^{-1-j}$.
For the remaining terms, use the bound
$(1+r_z/\sqrt t)^{-m}\le(\sqrt t/r_z)^m$ and sum
$t^{(m-3-j)/2}$ up to $r_z^2$; the same bound results.
This proves \eqref{eq:heatsum} and the asserted pointwise convergence
of $v$.
It gives the value and difference estimates in
\eqref{eq:smoothbounds}; for time derivatives use
\[
 \beta^{-1}(D_*+r_z)^{-3}\le W(z)/D_*.
\]

For the point part, \eqref{eq:heatderivative} gives
\eqref{eq:pointbounds} with $m=12$.  Formula
\eqref{eq:slotdecomposition} adds only a site-supported slot deviation
decaying as $e^{-3\ell s}$.  This deviation, its spatial differences,
and its time derivative satisfy the same bounds, with constants
depending on $S$.  In the reverse direction the initial seam removes
that deviation; at a seam its change is bounded by the incoming
absolute value of the slot deviation and hence by $C W_d(z)$.
The smooth part is unchanged
because it is slot-uniform.

It remains to justify multiplication by $\chi$.  Value and time
derivative estimates are preserved.  In a spatial difference the
extra term contains a difference of $\chi$, supported where
$r_z$ is of order $R$.  There
$G(z)/R\le C W(z)$, and
$P_d(z)/R\le C W_d(z)$ since $d\le\sqrt{1+\beta}\ll R$.
Thus the difference estimates are preserved as well.
Finally, if $s\ge a\beta$, then $d$ is comparable with $D_*$;
direct comparison of the envelopes gives
$P_d\le C_aG$, $W_d\le C_aW$, and
$W_d/d\le C_aW/D_*$.  This proves the last assertion.
\end{proof}

At the source site, the free potential therefore recovers the ideal
density exactly.  Fourier inversion and monotone convergence give
\begin{equation}\label{eq:freeidealdensity}
 \begin{split}
 \ell v(x)=\sum_{n\ge1}p_{n\beta}(0)
 &=\int_{[-\pi,\pi]^3}\frac{d^3k}{(2\pi)^3}
      \sum_{n\ge1}e^{-2n\beta S\sum_j(1-\cos k_j)}\\
 &=n_S^{\rm sw}(\beta).
 \end{split}
\end{equation}
The integral is finite in dimension three, also as a consequence of
the heat-kernel sum in \eqref{eq:heatsum}.

The uncut potential need not be square summable.  Its evolutions
above are defined pointwise by the convergent positive series, while
the trial vector and every restricted profile used in a finite cube
have finite support.
Equation~\eqref{eq:freeidealdensity} identifies $v(x)$ with the
target in Proposition~\ref{prop:comparison}.  It remains to bound
the change of return capacity when the ideal period is replaced by
the walk among holes.
Define $\mu_+,\mu_-$ from this $h$ and $c$ by
\eqref{eq:capacityresiduals}.  Section~\ref{sec:sources} realizes
these residuals as propagated local errors of the profiles.
